\ProvidesFile{algebra.tex}[2026/07/26 Algebraic and semilattice notation]

% Algebraische Halbverbands- und Verbandsprädikate.
\DeclareRobustCommand{\BinOp}[1]{%
  \ensuremath{\mathbin{#1}}%
}
\DeclareRobustCommand{\JoinOp}{%
  \ensuremath{\mathbin{\sqcup}}%
}
\DeclareRobustCommand{\MeetOp}{%
  \ensuremath{\mathbin{\sqcap}}%
}
\DeclareRobustCommand{\FamilyJoinOp}[1]{%
  \ensuremath{\mathbin{\sqcup_{#1}}}%
}
\DeclareRobustCommand{\NeutralSemiOp}{%
  \ensuremath{\mathbin{\star}}%
}
% Quelltextmarke fuer eine kontextuell festgelegte, in der gesetzten Formel
% jedoch tatsaechlich weggelassene Halbverbandsoperation. Abstrakte Produkte
% erscheinen damit als Juxtaposition xy; konkrete Operatoren bleiben benannt.
\DeclareRobustCommand{\OmittedSemiOp}{}
\DeclareRobustCommand{\SemigroupAlg}[2]{%
  \ensuremath{\mathsf{HGr}\!\left(#1,#2\right)}%
}
\DeclareRobustCommand{\CommutativeSemigroupAlg}[2]{%
  \ensuremath{\mathsf{KHGr}\!\left(#1,#2\right)}%
}
\DeclareRobustCommand{\MonoidAlg}[3]{%
  \ensuremath{\mathsf{Mon}\!\left(#1,#2,#3\right)}%
}
\DeclareRobustCommand{\CommutativeMonoidAlg}[3]{%
  \ensuremath{\mathsf{KMon}\!\left(#1,#2,#3\right)}%
}
\DeclareRobustCommand{\GroupAlg}[3]{%
  \ensuremath{\mathsf{Grp}\!\left(#1,#2,#3\right)}%
}
\DeclareRobustCommand{\AbelianGroupAlg}[3]{%
  \ensuremath{\mathsf{AGrp}\!\left(#1,#2,#3\right)}%
}
% Five-place predicates use the order
% (carrier, addition, multiplication, additive zero, multiplicative one).
\DeclareRobustCommand{\SemiringAlg}[5]{%
  \ensuremath{\mathsf{Hring}\!\left(#1,#2,#3,#4,#5\right)}%
}
\DeclareRobustCommand{\CommutativeSemiringAlg}[5]{%
  \ensuremath{\mathsf{KHring}\!\left(#1,#2,#3,#4,#5\right)}%
}
\DeclareRobustCommand{\RingAlg}[5]{%
  \ensuremath{\mathsf{Ring}\!\left(#1,#2,#3,#4,#5\right)}%
}
\DeclareRobustCommand{\CommutativeRingAlg}[5]{%
  \ensuremath{\mathsf{KRing}\!\left(#1,#2,#3,#4,#5\right)}%
}

% Homomorphismen und Isomorphismen algebraischer Strukturen.
% Die beiden Strukturargumente sind bewusst jeweils eine geschweifte
% Kommaliste. Dadurch bleiben auch die zweioperativen Strukturen vollständig
% lesbar, ohne die TeX-Grenze von neun Makroargumenten zu überschreiten:
%   \SemigroupHom{f}{A,\star}{B,\diamond}
%   \RingHom{f}{R,+,\cdot,0_R,1_R}{S,\oplus,\odot,0_S,1_S}.
\DeclareRobustCommand{\SemigroupHom}[3]{%
  \ensuremath{\mathsf{HGrHom}\!\left(#1;#2;#3\right)}%
}
\DeclareRobustCommand{\SemigroupIso}[3]{%
  \ensuremath{\mathsf{HGrIso}\!\left(#1;#2;#3\right)}%
}
\DeclareRobustCommand{\CommutativeSemigroupHom}[3]{%
  \ensuremath{\mathsf{KHGrHom}\!\left(#1;#2;#3\right)}%
}
\DeclareRobustCommand{\CommutativeSemigroupIso}[3]{%
  \ensuremath{\mathsf{KHGrIso}\!\left(#1;#2;#3\right)}%
}
\DeclareRobustCommand{\IdempotentSemigroupHom}[3]{%
  \ensuremath{\mathsf{IdemHGrHom}\!\left(#1;#2;#3\right)}%
}
\DeclareRobustCommand{\IdempotentSemigroupIso}[3]{%
  \ensuremath{\mathsf{IdemHGrIso}\!\left(#1;#2;#3\right)}%
}
\DeclareRobustCommand{\CommutativeIdempotentSemigroupHom}[3]{%
  \ensuremath{\mathsf{KIHGrHom}\!\left(#1;#2;#3\right)}%
}
\DeclareRobustCommand{\CommutativeIdempotentSemigroupIso}[3]{%
  \ensuremath{\mathsf{KIHGrIso}\!\left(#1;#2;#3\right)}%
}
\DeclareRobustCommand{\CancellativeSemigroupHom}[3]{%
  \ensuremath{\mathsf{K\ddot{u}rzHGrHom}\!\left(#1;#2;#3\right)}%
}
\DeclareRobustCommand{\CancellativeSemigroupIso}[3]{%
  \ensuremath{\mathsf{K\ddot{u}rzHGrIso}\!\left(#1;#2;#3\right)}%
}
\DeclareRobustCommand{\MonoidHom}[3]{%
  \ensuremath{\mathsf{MonHom}\!\left(#1;#2;#3\right)}%
}
\DeclareRobustCommand{\MonoidIso}[3]{%
  \ensuremath{\mathsf{MonIso}\!\left(#1;#2;#3\right)}%
}
\DeclareRobustCommand{\CommutativeMonoidHom}[3]{%
  \ensuremath{\mathsf{KMonHom}\!\left(#1;#2;#3\right)}%
}
\DeclareRobustCommand{\CommutativeMonoidIso}[3]{%
  \ensuremath{\mathsf{KMonIso}\!\left(#1;#2;#3\right)}%
}
\DeclareRobustCommand{\SemiringHom}[3]{%
  \ensuremath{\mathsf{HringHom}\!\left(#1;#2;#3\right)}%
}
\DeclareRobustCommand{\SemiringIso}[3]{%
  \ensuremath{\mathsf{HringIso}\!\left(#1;#2;#3\right)}%
}
\DeclareRobustCommand{\CommutativeSemiringHom}[3]{%
  \ensuremath{\mathsf{KHringHom}\!\left(#1;#2;#3\right)}%
}
\DeclareRobustCommand{\CommutativeSemiringIso}[3]{%
  \ensuremath{\mathsf{KHringIso}\!\left(#1;#2;#3\right)}%
}
\DeclareRobustCommand{\GroupHom}[3]{%
  \ensuremath{\mathsf{GrpHom}\!\left(#1;#2;#3\right)}%
}
\DeclareRobustCommand{\GroupIso}[3]{%
  \ensuremath{\mathsf{GrpIso}\!\left(#1;#2;#3\right)}%
}
\DeclareRobustCommand{\AbelianGroupHom}[3]{%
  \ensuremath{\mathsf{AGrpHom}\!\left(#1;#2;#3\right)}%
}
\DeclareRobustCommand{\AbelianGroupIso}[3]{%
  \ensuremath{\mathsf{AGrpIso}\!\left(#1;#2;#3\right)}%
}
\DeclareRobustCommand{\RingHom}[3]{%
  \ensuremath{\mathsf{RingHom}\!\left(#1;#2;#3\right)}%
}
\DeclareRobustCommand{\RingIso}[3]{%
  \ensuremath{\mathsf{RingIso}\!\left(#1;#2;#3\right)}%
}
\DeclareRobustCommand{\CommutativeRingHom}[3]{%
  \ensuremath{\mathsf{KRingHom}\!\left(#1;#2;#3\right)}%
}
\DeclareRobustCommand{\CommutativeRingIso}[3]{%
  \ensuremath{\mathsf{KRingIso}\!\left(#1;#2;#3\right)}%
}
\DeclareRobustCommand{\SemilatticeHom}[1]{%
  \ensuremath{\mathsf{HVHom}\!\left(#1\right)}%
}
\DeclareRobustCommand{\SemilatticeIso}[1]{%
  \ensuremath{\mathsf{HVIso}\!\left(#1\right)}%
}
\DeclareRobustCommand{\LatticeHom}[1]{%
  \ensuremath{\mathsf{VerHom}\!\left(#1\right)}%
}
\DeclareRobustCommand{\LatticeIso}[1]{%
  \ensuremath{\mathsf{VerIso}\!\left(#1\right)}%
}
\DeclareRobustCommand{\SemiringSuccessor}[1]{%
  \ensuremath{\sigma_{#1}}%
}
\DeclareRobustCommand{\NaturalSemiringMap}[1]{%
  \ensuremath{\eta_{#1}}%
}
\DeclareRobustCommand{\IntegerRingMap}[1]{%
  \ensuremath{\zeta_{#1}}%
}
\DeclareRobustCommand{\IntegerRingPairValue}[3]{%
  \ensuremath{\vartheta_{#1}\!\left(#2,#3\right)}%
}
\DeclareRobustCommand{\InducedLeq}{%
  \ensuremath{\mathrel{\leq}}%
}
\DeclareRobustCommand{\InducedGeq}{%
  \ensuremath{\mathrel{\geq}}%
}
\DeclareRobustCommand{\Semi}[1]{%
  \ensuremath{\mathsf{HV}\!\left(#1\right)}%
}
\DeclareRobustCommand{\LatticeAlg}[1]{%
  \ensuremath{\mathsf{Verband}\!\left(#1\right)}%
}

% Universelle Eigenschaften zweier Elemente.
\DeclareRobustCommand{\PairSup}[5]{%
  \ensuremath{#5=\sup_{#1,#2}\!\left(\{#3,#4\}\right)}%
}
\DeclareRobustCommand{\PairInf}[5]{%
  \ensuremath{#5=\inf_{#1,#2}\!\left(\{#3,#4\}\right)}%
}
\DeclareRobustCommand{\PairJoinOp}[1]{%
  \ensuremath{\mathsf{SupOp}\!\left(#1\right)}%
}

% Supremums-Halbverbände mit Null, beschränkte Supremums-Halbverbände
% mit Null und Eins sowie Irreduzibilität.
\DeclareRobustCommand{\ZeroJoinSemi}[1]{%
  \ensuremath{\mathsf{SupHV}_{0}\!\left(#1\right)}%
}
\DeclareRobustCommand{\BoundedJoinSemi}[1]{%
  \ensuremath{\mathsf{SupHV}_{0,1}\!\left(#1\right)}%
}
\DeclareRobustCommand{\ZeroJoinSemiHom}[1]{%
  \ensuremath{\mathsf{SupHV}_{0}\mathsf{Hom}\!\left(#1\right)}%
}
\DeclareRobustCommand{\ZeroJoinSemiIso}[1]{%
  \ensuremath{\mathsf{SupHV}_{0}\mathsf{Iso}\!\left(#1\right)}%
}
\DeclareRobustCommand{\BoundedJoinSemiHom}[1]{%
  \ensuremath{\mathsf{SupHV}_{0,1}\mathsf{Hom}\!\left(#1\right)}%
}
\DeclareRobustCommand{\BoundedJoinSemiIso}[1]{%
  \ensuremath{\mathsf{SupHV}_{0,1}\mathsf{Iso}\!\left(#1\right)}%
}
\DeclareRobustCommand{\JoinIrr}[2]{%
  \ensuremath{\mathsf{Irr}_{#1}\!\left(#2\right)}%
}
\DeclareRobustCommand{\JoinIrrSet}[1]{%
  \ensuremath{\mathop{\mathrm{Irr}}\nolimits_{#1}}%
}

% Hauptfilter und Hauptideal in einer dargestellten Ordnung.
\DeclareRobustCommand{\PrincipalFilter}[2]{%
  \ensuremath{\mathord{\uparrow}_{#1}\!#2}%
}
\DeclareRobustCommand{\PrincipalIdeal}[2]{%
  \ensuremath{\mathord{\downarrow}_{#1}\!#2}%
}

% Injektive Form von "höchstens die Hälfte" und Frankls Halbverbandsform.
\DeclareRobustCommand{\AtMostHalf}[2]{%
  \ensuremath{\mathsf{Anteil}_{\leq 1/2}\!\left(#1;#2\right)}%
}
\DeclareRobustCommand{\RareOcc}[2]{%
  \ensuremath{\mathsf{Selten}\!\left(#1;#2\right)}%
}
\DeclareRobustCommand{\FranklSemi}[1]{%
  \ensuremath{\mathsf{FranklHV}\!\left(#1\right)}%
}
\DeclareRobustCommand{\InterClosedFam}[1]{%
  \ensuremath{\mathsf{SchnittAbg}\!\left(#1\right)}%
}

% Kanonische Familie der unter einem Element liegenden \sqcup-Irreduziblen.
\DeclareRobustCommand{\IrrBelow}[2]{%
  \ensuremath{\mathsf{IrrProfil}_{#1}\!\left(#2\right)}%
}
\DeclareRobustCommand{\CanonicalMeetFamily}[1]{%
  \ensuremath{\mathsf{ProfilFamilie}\!\left(#1\right)}%
}
\DeclareRobustCommand{\CanonicalProfileMap}[1]{%
  \ensuremath{c_{#1}}%
}
\DeclareRobustCommand{\JoinSeparatorSet}[4]{%
  \ensuremath{\mathsf{Trenner}_{#1,#2}\!\left(#3,#4\right)}%
}
\DeclareRobustCommand{\ProperLowerSet}[3]{%
  \ensuremath{\mathord{\downarrow}^{-}_{#1,#2}\!#3}%
}
\DeclareRobustCommand{\FranklComplementMap}[1]{%
  \ensuremath{\kappa_{#1}}%
}

% Kurzzeichen der globalen Formen von Frankls Vermutung. Die Mengen- und die
% Halbverbandsform erhalten in B18 ihre vollständigen registrierten Definitionen;
% die Kurzzeichen halten insbesondere Lemmon-Tabellen lesbar.
\DeclareRobustCommand{\FranklSetStatement}{\ensuremath{\mathsf{FV}_{\mathrm{Menge}}}}
\DeclareRobustCommand{\FranklInterStatement}{\ensuremath{\mathsf{FV}_{\cap}}}
\DeclareRobustCommand{\FranklJoinStatement}{\ensuremath{\mathsf{FV}_{\sqcup}}}
